%
% 14-prinzip.tex -- Die der Methode zugrunde liegenden Prinzipien
%
% (c) 2026 Prof Dr Andreas Müller
%
\subsection{Die der Methode zugrunde liegenden Prinzipien}
Es wurde bereits früher festgehalten, dass die dargestellte Methode
für ein Computeralgebrasystem nicht durchführbar ist, weil eine
Methode zur Bestimmung der Faktorisierung des Nenners in Linearfaktoren
oder quadratische reelle Faktoren fehlt.
\index{Computeralgebrasystem}%
Es muss daher eine alternative Vorgehensweise gefunden werden,
die mit einer leichter zu findenden Faktorisierung des Nenners
arbeiten kann.
Die Faktorisierung muss immer noch einige Eigenschaften erfüllen,
damit die einzelnen Schritte verallgemeinert werden können.

\begin{enumerate}
\item
Die Partialbruchzerlegung wurde dazu verwendet, die 
rationale Funktion in Term mit speziell einfachen Nennern
aufzuspalten.
\index{Partialbruchzerlegung}%
In jedem Term kamen höchstens Potenzen eines einzigen Faktors
des Nenners vor.
Wie in der Darstellung der Partialbruchzerlegung angedeutet
wurde, müssen die Faktoren teilerfremd sein.
\item
Für den Reduktionsschritt von
\index{Reduktion}%
Satz~\ref{buch:ratint:bernoulli:satz:reduktion}
wurde benötigt, dass sich jeder denkbare Zähler in der Form
\[
s(z) \cdot q(z) + t(z)\cdot q'(z)
\]
ausdrücken lässt, wobei $q(z)$ ein einzelner Faktor in der
Nennerfaktorisierung ist.
Dies ist genau dann möglich, wenn $q(z)$ und $q'(z)$ teilerfremd
sind, also keine gemeinsame Nullstelle haben.
\index{teilerfremd}%
Dies wieder ist gleichbedeutend damit, dass $q(z)$ nur einfache
Nullstellen hat.
\item
Es wird ein Algorithmus benötigt, der die Polynom $s(z)$ und
$t(z)$ auf effiziente Art und Weise finden kann.
\item
Es müssen Integrale der Form 
\[
\int \frac{p(z)}{q(z)}\,dz
\]
gefunden werden, wobei $p(z)$ ein beliebiges Polynom ist.
\end{enumerate}

Die Aufgabe von Punkt~3, das Finden einer Darstellung eines
beliebigen Zahlerpolynoms Zwecks Reduktion des Exponenten im
Nenner wird durch den euklidischen Algorithmus ermöglicht,
\index{euklidischer Algorithmus}%
der im Abschnitt~\ref{buch:ratint:section:euklid} dargestellt
wird.

Eine geeignete Faktorisierung, die immer gefunden werden kann,
ist die sogenannte quadratfreie Faktorisierung, die in Abschnitt 
\ref{buch:ratint:section:hermite}
zusammen mit der Reduktion der Integrale von rationalen
Funktionen mit Potenzen des Nennerfaktors vorgestellt wird.


